\documentclass{article}
\usepackage[utf8]{inputenc}
\usepackage[spanish,mexico]{babel}
\usepackage[margin=2.3cm]{geometry}
\usepackage{graphicx}
\usepackage[export]{adjustbox}
\usepackage{caption}
\usepackage{subcaption}
\usepackage{fancyhdr}
\usepackage{lipsum}
\usepackage{float}
\usepackage{amsfonts, amsthm}
\usepackage{upgreek}
\usepackage{physics}
\usepackage{cancel}
\usepackage{amssymb, latexsym, amsmath}
\usepackage{amsthm}

% Notas al pie.
\renewcommand{\thefootnote}{\Roman{footnote}}

% Evitar sangrías
\setlength{\parindent}{0cm}

% Letter colors.
\usepackage[usenames,dvipsnames,svgnames,table]{xcolor}

% Código
\newcommand{\code}[1]{\textcolor{white!25!black}{\texttt{#1}}}
\usepackage{listings}

% Def. Dr. César.
\usepackage{tikz}
\usetikzlibrary{shapes,calc}
\tikzstyle{edge}=[shorten <=2pt, shorten >=2pt, >=stealth, line width=1.1pt]
\tikzstyle{blueE}=[shorten <=2pt, shorten >=2pt, >=stealth, line width=1.5pt, blue]
\tikzstyle{blackV}=[circle, fill=black, minimum size=6pt, inner sep=0pt, outer sep=0pt]
\tikzstyle{blueV}=[circle, fill=blue, draw, minimum size=6pt, line width=0.75pt, inner sep=0pt, outer sep=0pt]
\tikzstyle{redV}=[circle, fill=red, draw, minimum size=6pt, line width=0.75pt, inner sep=0pt, outer sep=0pt]
\tikzstyle{redSV}=[semicircle, fill=red, minimum size=3pt, inner sep=0pt, outer sep=0pt, rotate=225]
\tikzstyle{blueSV}=[semicircle, fill=blue, minimum size=3pt, inner sep=0pt, outer sep=0pt, rotate=225]
\tikzstyle{blackSV}=[semicircle, fill=black, minimum size=3pt, inner sep=0pt, outer sep=0pt, rotate=225]
\tikzstyle{vertex}=[circle, draw, minimum size=6pt, line width=0.75pt, inner sep=0pt, outer sep=0pt]

% Proof
\renewcommand*{\proofname}{\textbf{Soluci\'on:}}

% Margins
%\addtolength{\textheight}{3cm}

\usepackage{amsmath}
\usepackage{amsmath}
\usepackage{amssymb}
\usepackage{amsthm}
\usepackage{url}
\usepackage{hyperref}

%%Tablas con color
\usepackage{xcolor, colortbl}
\usepackage{array, multirow, multicol}
\cellcolor[modelo color]{color}

\pagestyle{fancyplain}

\title{Facultad de Ciencias, UNAM}
\date{\today}

\hypersetup{
  colorlinks=true,
  linkcolor=blue,
  filecolor=magenta,      
  urlcolor=cyan,
  pdftitle={Overleaf Example},
  pdfpagemode=FullScreen,
}

\begin{document}
\fancyhead[r]{ Computación Distribuida 2023-1}
\thispagestyle{empty}

\begin{figure}[ht]
  \minipage{0.76\textwidth}
  \includegraphics[width=4cm]{Logo_UNAM.png}
  \label{EscudoUNAM}
  \endminipage
  \minipage{0.32\textwidth}
  \includegraphics[height = 4.9cm ,width=4cm]{Logo_FC.png}
  \label{EscudoFC}
  \endminipage
\end{figure}

\begin{center}
  \vspace{0.8cm}
  \LARGE
  UNIVERSIDAD NACIONAL AUTÓNOMA DE MÉXICO 
  
  \vspace{0.7cm}
  \LARGE
  FACULTAD DE CIENCIAS
  
  \vspace{0.8 cm}	
  \Large
  \textbf{Tarea 4}

  \vspace{0.8 cm}
  \normalsize	
  INTEGRANTES \\
  \vspace{.2cm}
  \large
  \textbf{Torres Valencia Kevin Jair - \texttt{318331818}}\\
  \textbf{Aguilera Moreno Adrián - \texttt{421005200}}\\
  \textbf{Natalia Abigail Pérez Romero  - \texttt{31814426}}\\
  %\textbf{Nombre - \texttt{Número de cuenta}}
  
  \vspace{1 cm}
  \normalsize	
  PROFESOR \\
  \vspace{.2cm}
  \large
  \textbf{Miguel Ángel Piña Avelino}
  
  \vspace{1 cm}
  AYUDANTE \\
  \vspace{.2cm}
  \large
  \textbf{Pablo Gerardo González López}
  \vspace{1.3cm}
  
  \normalsize	
  ASIGNATURA \\
  \vspace{.2cm}
  \large
  \textbf{Computación Distribuida}
  
  \vspace{1 cm}
  \today
\end{center}

%%%%%%%%%%% Aquí va la solución al problema 1.
\newpage

\section{Relojes Vectoriales}

Un reloj vectorial es un vector de $n$ escalares
\end{document}
